\chapter{Background}

\writingbox{Hook and intro to the background. State that there are large simplifications for understanding and modelling purposes?}

TRY NOT TO MAKE THIS TOO LONG. The focus should be on your work, not what already exists. Clear and concise.

\section{Auditory Localisation}

\writingbox{Frame an intro to this so that everything else is contextualised, this is like the expanded problem specification}

\section{Biological Foundations of Echolocation}

HAVE A FUNCTIONAL APPROACH TO THE BIOLOGY BACKGROUND. The accuracy is irrelevant, it is not your work, it is what you used to frame your model. Maybe explain this approach in the hook as well.

\subsection{Bats and Echolocation: The Signal}

Background on bats, and echolocation in general.

State the types of Echolocation, how it works, the types of bat calls, etc.

\subsection{Biological Neurons: The Basics of Biological Computation}

Maybe an explanation of biological neurons, inhibitory signals, etc, just so the rest of the biology makes more sense.

\subsection{Neurobiology of the Auditory Pathways}

Hook. Neurobiology of the auditory pathways of Microchiroptera

Figure of overall pathways

State that this will be the basis of the model developed.

Potentially state the current models of the biology here too? Maybe save it for later for ease of segmentation

\subsection{The Auditory Frontend: Peripheral \& Hindbrain Processing}

\subsubsection{Cochlea}

Biology of the cochlea. Brief explanation of the simplifications. Cochlear Nerve

\subsubsection{Cochlear Nuclei}

Cochlear nuclei (CN) intro: the first neurological processing unit

\paragraph{Ventral Cochlear Nucleus}

VCN explanation of PVCN and AVCN: Anteroventral Cochlear Nucleus (AVCN), Posteroventral Cochlear Nucleus (PVCN)

\paragraph{Dorsal Cochlear Nucleus}

DCN explanation

Maybe include a figure of the CN head


\subsection{The Distance Pathway: Temporal Interval Arithmetic}
\label{sec:distance_pathway}
Picks up the clean timestamps from the VCN. Focuses on the midbrain (Inferior Colliculus) mechanics: delay-tuned neurons, asymmetric sequential facilitation, and the Simmons-Suga hyperacuity debate.


\subsection{The Azimuth Pathway: Binaural Disparity Processing}
\label{sec:azimuth_pathway}
Picks up the clean timestamps from the VCN across both ears. Focuses on Interaural Time Differences (ITD) and Interaural Level Differences (ILD) lateral integration.

\subsection{The Elevation Pathway: Monaural Spectral Mapping}
\label{sec:elevation_pathway}
Picks up the output of the DCN and traces it to the collicular layers for vertical coordinate mapping.


\section{Neuromorphic Computing}

Pivot to how we model and compute using biological inspirations

\subsection{Computing with Spikes}

\subsubsection{Spikes and Encoding}

\subsubsection{Neurons}

LIF, RF, other types

\subsubsection{Neuronal Structures}

eg coincidence detection, comparison of LIF coincidence to cross correlation, spikes as samples and monte carlo ideas with bayesian brain theory, comparison of RF bank as FFT, switches idea of logic gates aetc

\subsection{Continuous Attractor Neural Networks}

Line attractor and ring model

\subsection{Spiking Neural Networks and Deep Learning}

Trainable systems, eg surrogate gradients, snntorch

\subsection{Neuromorphic Hardware}

Loihi, how it works and how we would have to design a simulation so that it is compatible

\section{Current Biomimetic Models of Echolocatory Bats}

Lay down what has already been done. 

THIS IS IMPORTANT. IT HIGHLIGHTS THE GAPS IN THE RESEARCH AND SHOWS THE ORIGINALITY OF YOUR MODEL.













